\documentclass[10pt,a4paper]{amsart}
\usepackage{amsmath,amssymb,amsthm,graphicx}
% Retained minimal preamble for the Stacks Project chapter flat.tex (More on Flatness).
% Verbatim excerpt of the chapter preamble: Xy-pic setup (native preamble lines 14-18) and the native
% theorem declarations with the mathematical macros (native preamble lines 154-217). Copied mechanically;
% nothing added, removed or redefined. The wrapper class and the amsmath/amssymb/amsthm/graphicx base
% packages are supplied by the assembler, so the class/externaldocument part of the native preamble is omitted.
% For commutative diagrams we use Xy-pic
\usepackage[all]{xy}

% We use 2cell for 2-commutative diagrams.
\xyoption{2cell}
% Theorem environments.
%
\theoremstyle{plain}
\newtheorem{theorem}[subsection]{Theorem}
\newtheorem{proposition}[subsection]{Proposition}
\newtheorem{lemma}[subsection]{Lemma}

\theoremstyle{definition}
\newtheorem{definition}[subsection]{Definition}
\newtheorem{example}[subsection]{Example}
\newtheorem{exercise}[subsection]{Exercise}
\newtheorem{situation}[subsection]{Situation}

\theoremstyle{remark}
\newtheorem{remark}[subsection]{Remark}
\newtheorem{remarks}[subsection]{Remarks}

\numberwithin{equation}{subsection}

% Macros
%
\def\lim{\mathop{\mathrm{lim}}\nolimits}
\def\colim{\mathop{\mathrm{colim}}\nolimits}
\def\Spec{\mathop{\mathrm{Spec}}}
\def\Hom{\mathop{\mathrm{Hom}}\nolimits}
\def\Ext{\mathop{\mathrm{Ext}}\nolimits}
\def\SheafHom{\mathop{\mathcal{H}\!\mathit{om}}\nolimits}
\def\SheafExt{\mathop{\mathcal{E}\!\mathit{xt}}\nolimits}
\def\Sch{\mathit{Sch}}
\def\Mor{\mathop{\mathrm{Mor}}\nolimits}
\def\Ob{\mathop{\mathrm{Ob}}\nolimits}
\def\Sh{\mathop{\mathit{Sh}}\nolimits}
\def\NL{\mathop{N\!L}\nolimits}
\def\CH{\mathop{\mathrm{CH}}\nolimits}
\def\proetale{{pro\text{-}\acute{e}tale}}
\def\etale{{\acute{e}tale}}
\def\QCoh{\mathit{QCoh}}
\def\Ker{\mathop{\mathrm{Ker}}}
\def\Im{\mathop{\mathrm{Im}}}
\def\Coker{\mathop{\mathrm{Coker}}}
\def\Coim{\mathop{\mathrm{Coim}}}

% Boxtimes
%
\DeclareMathSymbol{\boxtimes}{\mathbin}{AMSa}{"02}

%
% Macros for moduli stacks/spaces
%
\def\QCohstack{\mathcal{QC}\!\mathit{oh}}
\def\Cohstack{\mathcal{C}\!\mathit{oh}}
\def\Spacesstack{\mathcal{S}\!\mathit{paces}}
\def\Quotfunctor{\mathrm{Quot}}
\def\Hilbfunctor{\mathrm{Hilb}}
\def\Curvesstack{\mathcal{C}\!\mathit{urves}}
\def\Polarizedstack{\mathcal{P}\!\mathit{olarized}}
\def\Complexesstack{\mathcal{C}\!\mathit{omplexes}}
% \Pic is the operator that assigns to X its picard group, usage \Pic(X)
% \Picardstack_{X/B} denotes the Picard stack of X over B
% \Picardfunctor_{X/B} denotes the Picard functor of X over B
\def\Pic{\mathop{\mathrm{Pic}}\nolimits}
\def\Picardstack{\mathcal{P}\!\mathit{ic}}
\def\Picardfunctor{\mathrm{Pic}}
\def\Deformationcategory{\mathcal{D}\!\mathit{ef}}

% --- appended 2026-10-12 (Chief): the project own definitions of the reference/slogan/history
% display environments, copied mechanically from the shared Stacks preamble lines 8-12
% (\usepackage{verbatim} plus the three \newenvironment lines, which the chapter preamble keeps
% verbatim from the project preamble). Nothing added, removed or redefined: these environments
% carry no mathematical content and the project itself typesets them as comments. Without these
% three lines a counted statement whose own text uses them fails with Environment slogan undefined.
\usepackage{verbatim}
\newenvironment{reference}{\comment}{\endcomment}
\newenvironment{slogan}{\comment}{\endcomment}
\newenvironment{history}{\comment}{\endcomment}
\usepackage{cleveref}
\crefname{originaldefinition}{original definition}{original definitions}
\Crefname{originaldefinition}{Original definition}{Original definitions}
\makeatletter
\newlabel{definition-h-covering}{{0ETS}{source}{}{sourcepointer-1}{}}
\newlabel{definition-h-covering@cref}{{[originaldefinition][0][]0ETS}{[1][0][]source}{}{}{}}
\makeatother
\crefname{originaldefinition}{original definition}{original definitions}
\Crefname{originaldefinition}{Original definition}{Original definitions}
\makeatletter
\newlabel{topologies-definition-ph-covering}{{0DBG}{source}{}{sourcepointer-2}{}}
\newlabel{topologies-definition-ph-covering@cref}{{[originaldefinition][0][]0DBG}{[1][0][]source}{}{}{}}
\makeatother
\crefname{originallemma}{original lemma}{original lemmas}
\Crefname{originallemma}{Original lemma}{Original lemmas}
\makeatletter
\newlabel{limits-lemma-proper-limit-of-proper-finite-presentation}{{09ZR}{source}{}{sourcepointer-3}{}}
\newlabel{limits-lemma-proper-limit-of-proper-finite-presentation@cref}{{[originallemma][0][]09ZR}{[1][0][]source}{}{}{}}
\makeatother
\crefname{originallemma}{original lemma}{original lemmas}
\Crefname{originallemma}{Original lemma}{Original lemmas}
\makeatletter
\newlabel{limits-lemma-descend-opens}{{01Z4}{source}{}{sourcepointer-4}{}}
\newlabel{limits-lemma-descend-opens@cref}{{[originallemma][0][]01Z4}{[1][0][]source}{}{}{}}
\makeatother
\crefname{originallemma}{original lemma}{original lemmas}
\Crefname{originallemma}{Original lemma}{Original lemmas}
\makeatletter
\newlabel{limits-lemma-limit-affine}{{01Z6}{source}{}{sourcepointer-5}{}}
\newlabel{limits-lemma-limit-affine@cref}{{[originallemma][0][]01Z6}{[1][0][]source}{}{}{}}
\makeatother
\crefname{originalproposition}{original proposition}{original propositions}
\Crefname{originalproposition}{Original proposition}{Original propositions}
\makeatletter
\newlabel{limits-proposition-characterize-locally-finite-presentation}{{01ZC}{source}{}{sourcepointer-6}{}}
\newlabel{limits-proposition-characterize-locally-finite-presentation@cref}{{[originalproposition][0][]01ZC}{[1][0][]source}{}{}{}}
\makeatother
\title{Every h cover of an affine scheme is refined by a finite affine cover of a surjective proper finitely presented morphism}
\author{The Stacks Project Authors}
\date{stacks-project:flat.tex:local-snapshot-20261011; GNU FDL 1.2 or later}
\begin{document}
\maketitle
\noindent\textit{Editorial note.} All delivered mathematical text is the author's own, copied from the pinned Stacks Project chapter file flat.tex; no statement or proof step was rewritten, reordered or removed.\par\medskip
\noindent\textit{Editorial note.} The preamble is the chapter's own retained minimal preamble (sources/stacks-flat/retained-preamble.tex), a verbatim excerpt of the native Xy-pic setup and of the native theorem declarations and mathematical macros, followed by the project own definitions of the reference, slogan and history display environments (verbatim package plus three newenvironment lines) that the retained file carries at its end and that the project itself typesets as comments. No macro was invented or redefined, and the wrapper class and the amsmath/amssymb/amsthm/graphicx base packages are supplied by the assembler.\par\medskip
\noindent\textit{Editorial note.} The counted statement and proof contain no cross reference, no citation and no included file, so the delivered body is self contained.\par\medskip
\noindent\textit{Editorial note.} Canonical location of the counted proposition: Stacks tag 0ETU (sources/stacks-flat/pinned-original-source.tex line 10303).\par\medskip
\noindent\textit{Allowed standard prerequisites (editorial).}\begin{itemize}
\item Definition of an h covering (and of the ph covering used in the comparison). Same chapter, Definition labelled definition-h-covering, cited at proof lines 10322-10325, with Topologies, Definition labelled topologies-definition-ph-covering. Supplies the initial proper surjective cover and its refining finite affine cover.
\item Approximating a proper morphism by proper finitely presented morphisms. Limits, Lemma labelled limits-lemma-proper-limit-of-proper-finite-presentation, cited at proof lines 10327-10329. Writes Y -> X as a directed limit of proper finitely presented morphisms Y\_lambda -> X with surjectivity preserved.
\item Descending opens and affineness through a directed limit. Limits, Lemmas labelled limits-lemma-descend-opens and limits-lemma-limit-affine, cited at proof lines 10335-10341. Spreads the finite affine cover to some Y\_lambda.
\item The functorial characterisation of local finite presentation. Limits, Proposition labelled limits-proposition-characterize-locally-finite-presentation, cited at proof lines 10343-10347. Descends the refinement morphisms h\_j to Y\_\{lambda,j\}.
\end{itemize}\medskip
\ifdefined\hypertarget\hypertarget{sourcepointer-1}{}\fi
\noindent\textit{Source reference.} Original definition 0ETS: sources/stacks-flat/pinned-original-source.tex, source lines 10171--10171. Same-chapter pointer into the author environment definition at source line 10171; invoked at proof line 10321 of the counted proof. It is retained only as a navigation pointer; every load-bearing step is inside the delivered span.\par\medskip
\ifdefined\hypertarget\hypertarget{sourcepointer-2}{}\fi
\noindent\textit{Source reference.} Original definition 0DBG: sources/stacks-chow/tags.txt, source lines 16696--16696. Cross-chapter standard prerequisite located in the project tag list staged in the local snapshot; invoked at proof line 10322 of the counted proof. It is retained only as a navigation pointer; every load-bearing step is inside the delivered span.\par\medskip
\ifdefined\hypertarget\hypertarget{sourcepointer-3}{}\fi
\noindent\textit{Source reference.} Original lemma 09ZR: sources/stacks-chow/tags.txt, source lines 12609--12609. Cross-chapter standard prerequisite located in the project tag list staged in the local snapshot; invoked at proof line 10326 of the counted proof. It is retained only as a navigation pointer; every load-bearing step is inside the delivered span.\par\medskip
\ifdefined\hypertarget\hypertarget{sourcepointer-4}{}\fi
\noindent\textit{Source reference.} Original lemma 01Z4: sources/stacks-chow/tags.txt, source lines 2579--2579. Cross-chapter standard prerequisite located in the project tag list staged in the local snapshot; invoked at proof line 10335 of the counted proof. It is retained only as a navigation pointer; every load-bearing step is inside the delivered span.\par\medskip
\ifdefined\hypertarget\hypertarget{sourcepointer-5}{}\fi
\noindent\textit{Source reference.} Original lemma 01Z6: sources/stacks-chow/tags.txt, source lines 2581--2581. Cross-chapter standard prerequisite located in the project tag list staged in the local snapshot; invoked at proof line 10342 of the counted proof. It is retained only as a navigation pointer; every load-bearing step is inside the delivered span.\par\medskip
\ifdefined\hypertarget\hypertarget{sourcepointer-6}{}\fi
\noindent\textit{Source reference.} Original proposition 01ZC: sources/stacks-chow/tags.txt, source lines 2587--2587. Cross-chapter standard prerequisite located in the project tag list staged in the local snapshot; invoked at proof line 10347 of the counted proof. It is retained only as a navigation pointer; every load-bearing step is inside the delivered span.\par\medskip
\begin{lemma}
\label{lemma-approximate-h-cover}
Let $X$ be an affine scheme. Let $\{X_i \to X\}_{i \in I}$
be an h covering. Then there exists a surjective proper morphism
$$
Y \longrightarrow X
$$
of finite presentation (!) and a finite affine open covering
$Y = \bigcup_{j = 1, \ldots, m} Y_j$ such that
$\{Y_j \to X\}_{j = 1, \ldots, m}$ refines $\{X_i \to X\}_{i \in I}$.
\end{lemma}

\begin{proof}
By assumption there exists a proper surjective morphism
$Y \to X$ and a finite affine open covering
$Y = \bigcup_{j = 1, \ldots, m} Y_j$ such that
$\{Y_j \to X\}_{j = 1, \ldots, m}$ refines $\{X_i \to X\}_{i \in I}$.
This means that for each $j$ there is an index $i_j \in I$
and a morphism $h_j : Y_j \to X_{i_j}$ over $X$.
See Definition \ref{definition-h-covering} and
Topologies, Definition \ref{topologies-definition-ph-covering}.
The problem is that we don't know that $Y \to X$ is of finite
presentation.
By
Limits, Lemma \ref{limits-lemma-proper-limit-of-proper-finite-presentation}
we can write
$$
Y = \lim Y_\lambda
$$
as a directed limit of schemes $Y_\lambda$ proper and of finite presentation
over $X$ such that the morphisms $Y \to Y_\lambda$ and the
the transition morphisms are closed immersions. Observe that
each $Y_\lambda \to X$ is surjective.
By Limits, Lemma \ref{limits-lemma-descend-opens}
we can find a $\lambda$ and quasi-compact opens
$Y_{\lambda, j} \subset Y_\lambda$, $j = 1, \ldots, m$
covering $Y_\lambda$ and restricting to $Y_j$ in $Y$.
Then $Y_j = \lim Y_{\lambda, j}$.
After increasing $\lambda$ we may assume $Y_{\lambda, j}$
is affine for all $j$, see
Limits, Lemma \ref{limits-lemma-limit-affine}.
Finally, since $X_i \to X$ is locally of finite presentation
we can use the functorial characterization of morphisms
which are locally of finite presentation
(Limits, Proposition
\ref{limits-proposition-characterize-locally-finite-presentation})
to find a $\lambda$ such that for each $j$ there is
a morphism $h_{\lambda, j} : Y_{\lambda, j} \to X_{i_j}$
whose restriction to $Y_j$ is the morphism $h_j$ chosen above.
Thus $\{Y_{\lambda, j} \to X\}$ refines
$\{X_i \to X\}$ and the proof is complete.
\end{proof}

\end{document}
